\scanpage{60}
$B'$ doit être vide, au dessus de chaque $I_\alpha$ non tordu, il est \deleted{\unsure{ou}} vide ou une des deux orbites de $B_\alpha$ sur $I_\alpha$ (trois cas), de sorte que le nb de \inserted{telles} facettes \deleted{\unsure{est}} (y inclus la facette $C$)
\[
\deleted{\prod_{\unsure{\alpha\ \text{non tordu}}}}3^{\sum\alpha_i^+}
\]

Les facettes de codim donnée $d$ correspondent aux parties $B'$ telles que \deleted{\unsure{l'image}} \inserted{$\operatorname{Card}B'=d$} \deleted{\unsure{est}}. On trouve, si $f_d$ est le nb de ces facettes
\[
\boxed{\sum_{i=0}^n f_d t^d=\prod_{i=1}^n(1+2t^i)^{\alpha_i^+}}
\]

\uline{Exemple} Automorphismes de $C_3$

\margin{$+,8$\\$-,8$}
a) Autom ``indécomposables''
\[
\left.
\begin{aligned}
1^\circ)\\2^\circ)
\end{aligned}
\right\}
\begin{aligned}
&\text{non tordu :}\quad e_1\unsure{-|\to}e_2-e_3\to e_1\quad\oplus\\
&\text{tordu}\qquad e_1-e_2-e_3\to -e_1\quad\ominus
\end{aligned}
\]
Le premier est une rotation \inserted{d'ordre 3} autour de la diagonale joignant $\deleted{\unsure{z}}e_1+e_2+e_3$ et $-e_1-e_2-e_3$. Il y a (sans compter la facette pleine) \uline{2 facettes} fixes, savoir les deux sommets précédents.

Le deuxième est d'ordre 6, \deleted{\unsure{Et}} c'est \inserted{(mod. conj.)} le produit du précédent avec l'antipodisme [\uline{NB} cette relation entre cas tordu et non tordu marche chaque fois que la dim $n$ est impaire]

b) Autom décomposables : deux cycles (dans $I$)

\deleted{$\unsure{\alpha_1^+,\alpha_1^+|\alpha_1^+\ \alpha_1^+}\ \unsure{f}\ \unsure{\alpha_2}$}
% archive sheet number: 59
